\chapter{Localidad de Bell y proyección holográfica}
\label{chap:localidad-bell-proyeccion}
\index{Bell, desigualdad de}\index{CHSH}\index{EPR}
\index{localidad}\index{no señalización}\index{independencia de ajustes}

La formulación EPR, el teorema de Bell y la condición de no señalización
responden a preguntas distintas \cite{EinsteinPodolskyRosen1935,Bell1964,
ClauserHorneShimonyHolt1969}. Una memoria global compartida no constituye por
sí sola una violación de localidad de Bell; tampoco una base ternaria o la
elección de \(\log 3\) como unidad de información modifican una desigualdad
de correlación. Esta sección fija los tipos necesarios y determina con
exactitud qué está probado y qué permanece abierto en HMT--MD.

\section{Escenario bipartito y premisas de localidad}
\label{sec:escenario-bell-hmt}

Considérense ajustes \(a,b\in\{0,1\}\), resultados
\(A,B\in\{-1,+1\}\) y una variable compartida
\(\lambda\in\Lambda\). Sea \(\rho\) una medida de probabilidad sobre
\(\Lambda\). La expresión
\[
 P(A,B\mid a,b)
 =\int_\Lambda
 P_A(A\mid a,\lambda)\,
 P_B(B\mid b,\lambda)\,d\rho(\lambda)
 \tag{B1}\label{eq:factorizacion-bell-local}
\]
reúne dos premisas diferentes:

\begin{enumerate}
\item \emph{factorización local}: condicionado por \(\lambda\), el resultado
de cada ala depende sólo de su ajuste local;
\item \emph{independencia de ajustes}: la medida satisface
\(d\rho(\lambda\mid a,b)=d\rho(\lambda)\).
\end{enumerate}

Se exige, además, normalización de las probabilidades y ausencia de una
selección de datos dependiente de \(a,b,A,B\). La terminología
\emph{localidad de Bell} se reservará aquí para este conjunto explícito de
hipótesis; no designa únicamente ausencia de señal superlumínica.

Definimos
\[
 E_{ab}:=\sum_{A,B=\pm1}AB\,P(A,B\mid a,b)
\]
y el funcional CHSH
\[
 S_{\rm CHSH}:=E_{00}+E_{01}+E_{10}-E_{11}.
 \tag{B2}\label{eq:funcional-chsh}
\]

\begin{theorem}[Cota CHSH bajo factorización local]
\label{thm:cota-chsh-local}
Si se satisfacen \eqref{eq:factorizacion-bell-local}, la independencia de
ajustes, la normalización y la ausencia de postselección dependiente de los
datos, entonces
\[
 \bigl|S_{\rm CHSH}\bigr|\leq2.
 \tag{B3}\label{eq:cota-chsh-dos}
\]
\end{theorem}

\begin{proof}
Para cada \(\lambda\), sean
\[
 x_a(\lambda):=\sum_{A=\pm1}A\,P_A(A\mid a,\lambda),
 \qquad
 y_b(\lambda):=\sum_{B=\pm1}B\,P_B(B\mid b,\lambda).
\]
Se tiene \(|x_a|,|y_b|\leq1\). La contribución condicionada al funcional es
\[
 s(\lambda)
 =x_0(y_0+y_1)+x_1(y_0-y_1).
\]
Por la desigualdad triangular,
\[
 |s(\lambda)|
 \leq |y_0+y_1|+|y_0-y_1|\leq2,
\]
pues \(y_0,y_1\in[-1,1]\). Integrar respecto de la misma medida \(\rho\),
independiente de \(a,b\), conserva la cota.
\end{proof}

\begin{corollary}[No señalización y su recíproco falso]
\label{cor:no-senalizacion-bell}
Las hipótesis del teorema implican
\[
 \sum_B P(A,B\mid a,b)=P(A\mid a),
 \qquad
 \sum_A P(A,B\mid a,b)=P(B\mid b).
 \tag{B4}\label{eq:no-senalizacion-bell}
\]
La no señalización, considerada aisladamente, no implica la factorización
local de \eqref{eq:factorizacion-bell-local}.
\end{corollary}

\begin{proof}
La primera identidad se obtiene sumando el factor \(P_B\), cuya
normalización vale uno; la segunda es simétrica. La implicación recíproca no
forma parte del argumento y existen distribuciones no señalizantes fuera del
politopo local.
\end{proof}

\begin{statusbox}
El teorema y el corolario son resultados
\textsc{exactos internos} de teoría de probabilidades finita. Su aplicación
a un aparato HMT--MD está
\textsc{demostrada con estructura de partida explícita} cuando se construyen
la medida, los ajustes y los mapas de resultado.
\end{statusbox}

\section{Proyección Bell del estado enriquecido común y preservación del estatuto físico}
\label{sec:estado-enriquecido-bell}

Sea \(X_{\rm enr}\) un espacio de estados enriquecidos y sean
\[
 r_A:\{0,1\}\times X_{\rm enr}\longrightarrow[-1,1],
 \qquad
 r_B:\{0,1\}\times X_{\rm enr}\longrightarrow[-1,1]
\]
dos lectores locales. Un estado puede contener memoria de frontera,
orientación, ruta y registro; toda esa información puede incorporarse a
\(\lambda=x\in X_{\rm enr}\).

\begin{theorem}[Una proyección local de un estado global sigue siendo
Bell-local]
\label{thm:no-promocion-bell-por-proyeccion}
Supóngase que:

\begin{enumerate}
\item \(x\in X_{\rm enr}\) se distribuye según una medida \(\mu\)
independiente de \(a,b\);
\item el resultado de Alice depende sólo de \((a,x)\) y el de Bob sólo de
\((b,x)\);
\item el promedio usa todas las realizaciones sin selección contextual.
\end{enumerate}

Entonces las correlaciones inducidas satisfacen
\[
 E_{ab}=\int_{X_{\rm enr}}r_A(a,x)r_B(b,x)\,d\mu(x),
\qquad
 |S_{\rm CHSH}|\leq2.
 \tag{B5}\label{eq:no-go-proyeccion-local}
\]
\end{theorem}

\begin{proof}
Se aplica el argumento del \cref{thm:cota-chsh-local} con
\(\lambda=x\), \(x_a=r_A(a,x)\) e \(y_b=r_B(b,x)\). La riqueza interna de
\(x\) no modifica la cota mientras se conserven las tres hipótesis.
\end{proof}

Este resultado impide una inferencia improcedente: una extensión global con
memoria puede explicar por qué dos lecturas comparten estructura, pero no
produce por sí sola una violación de Bell. Para obtener
\(|S_{\rm CHSH}|>2\), una realización debe abandonar o reformular al menos una
de las premisas: factorización, independencia de ajustes o muestreo no
contextual. Si la proyección depende conjuntamente de \(a\) y \(b\), deja de
ser local en el sentido de \eqref{eq:factorizacion-bell-local}; todavía habría
que probar separadamente \eqref{eq:no-senalizacion-bell}.

\section{Distinción entre localidad, funcional ternario y dependencia contextual}
\label{sec:cotejo-bell-historico}

Intervienen tres construcciones matemáticas que no deben identificarse:

\begin{enumerate}
\item la frase arquitectónica según la cual una extensión conserva memoria
global y sus lecturas locales publican proyecciones parciales;
\item el funcional ternario cíclico \(\mathcal B_3\), cuyos resultados
admiten también el valor \(0\);
\item el funcional contextual de información trítica \(I_3^{\rm ctx}\),
para el cual puede ocurrir \(I_3^{\rm ctx}(\lambda;a,b)>0\).
\end{enumerate}

El primer estrato es una interpretación holográfica compatible con el estado
enriquecido. No especifica aún una distribución bipartita
\(P_{\rm HMT}(A,B\mid a,b)\). El segundo no es el funcional CHSH
\eqref{eq:funcional-chsh}: cambia el alfabeto de resultados, el muestreo y la
normalización. La pretendida anulación por ortogonalidad no es una identidad:
de que dos ejes sean ortogonales no se sigue que una de las dos proyecciones
de un vector se anule. En efecto, para
\(v=(1,1)\), \(u(0)=(1,0)\) y
\(u(\pi/2)=(0,1)\), ambas proyecciones valen \(1\). Por consiguiente, esa
anulación no puede utilizarse para acotar \(\mathcal B_3\).

La tercera construcción relaja expresamente la independencia de ajustes. Si
\(\kappa\leq pH\), con \(H>0\), se deduce
\(p\geq\kappa/H\), no \(p\leq\kappa/H\); la desigualdad no proporciona la
cota superior de \(p\) que requeriría la conclusión propuesta para CHSH.
Asimismo, una partición binaria admisible requiere un censo exhaustivo de
políticas; sin ese censo no se obtiene la cota. Por estas dos razones, la
desigualdad informacional \(S\leq3\) no constituye una cota de Bell local bajo
las premisas estándar.

La base del logaritmo sólo cambia la unidad:
\[
 I_3(X;Y)=\frac{I_e(X;Y)}{\log 3}.
 \tag{B6}\label{eq:cambio-base-informacion-bell}
\]
En particular, ni un trit ni \(\log 3\) seleccionan una distribución, prueban
factorización o alteran la cota \(|S_{\rm CHSH}|\leq2\).
El factor \(\log3\) de la escala térmica por decisión triádica de
\cref{def:escala-termica-triadica} es una formalización distinta: cuantifica
la información de tres alternativas equiprobables y no constituye una prueba
de Bell.

Ningún rótulo algorítmico sustituye un mapa con dominio, codominio y regla
capaz de producir correlaciones CHSH. La fórmula interpretativa «ER=EPR»
requiere un entrelazador explícito entre corte geométrico, estado bipartito,
álgebra de observables y probabilidades.

La construcción global--local demuestra la existencia del estado enriquecido
y sus proyecciones. Una realización de correlaciones de Bell debe añadir un
estado bipartito, dos familias de observables locales y una medida conjunta
que reproduzca las probabilidades. Sin esos datos no se deducen una violación
CHSH, la saturación de Tsirelson ni una equivalencia ER--EPR.

\section{Criterio de realización y protocolo de contraste}
\label{sec:programa-bell-hmt}

Una promoción científica exige construir, en este orden,
\[
\begin{aligned}
X_{\rm enr}
&\longrightarrow
\bigl(\Omega_A,\Omega_B,\mu\bigr)\\
&\longrightarrow
\bigl(r_A,r_B,\mathcal A,\mathcal B\bigr)\\
&\longrightarrow
P_{\rm HMT}(A,B\mid a,b)
\longrightarrow
\bigl(E_{00},E_{01},E_{10},E_{11},S_{\rm CHSH}\bigr),
\end{aligned}
\tag{B7}\label{eq:programa-bell-hmt}
\]
donde \(\mathcal A,\mathcal B\) son las reglas de elección de ajustes. Antes
de abrir resultados experimentales deben congelarse:

\begin{enumerate}
\item la medida \(\mu\) y el dominio de estados;
\item la regla que asigna ajustes y la hipótesis, o renuncia explícita, de
independencia;
\item los mapas de resultado y el tratamiento de salidas nulas;
\item la regla de muestreo, pérdidas y coincidencias;
\item las cuatro correlaciones y la convención de signos;
\item la prueba de no señalización;
\item el intervalo de confianza y el falsador previamente registrado.
\end{enumerate}

Sólo después puede evaluarse si el mecanismo conserva la localidad de la
dinámica interna y, a la vez, produce una distribución proyectada no
Bell-local. Esta doble afirmación no se deduce del lenguaje holográfico: debe
probarse sobre \eqref{eq:programa-bell-hmt}.

\begin{criterion}[Control de Bell, localidad y no localidad]
\label{crit:bell-localidad-no-localidad}
La realización queda refutada o indebidamente tipada si ocurre cualquiera de
los siguientes hechos:

\begin{enumerate}
\item se afirma una desigualdad sin definir \(P(A,B\mid a,b)\), los ajustes y
los resultados;
\item se conserva \eqref{eq:factorizacion-bell-local} con independencia de
ajustes y, no obstante, se certifica \(|S_{\rm CHSH}|>2\);
\item se usa \(I_3(\lambda;a,b)>0\) y se denomina al modelo
«independiente de ajustes»;
\item una postselección dependiente de ajustes o resultados se oculta en la
proyección;
\item la distribución proyectada permite señalización;
\item se identifica \(\mathcal B_3\), \(\log 3\) o «ER=EPR» con CHSH sin
construir el entrelazador correspondiente.
\end{enumerate}
\end{criterion}

El censo ejecutable asociado enumera las dieciséis estrategias deterministas
locales binarias y verifica que todas satisfacen
\(|S_{\rm CHSH}|=2\). Por convexidad, este control cubre el politopo local
binario completo; no constituye una simulación de una interfaz HMT aún
ausente.
